\section{Supports of general seminorms}\label{sec:general-supports}

Let $A$ be any commutative ring. These support statements apply to the absolute
seminorm spectrum without choosing a coefficient ring. For a relative spectrum,
the prescribed base seminorm forces a specific contraction of every support.

\begin{definition}[The support map]\label{spectrum:support}
\lean{EnvelopeProperty.Berkovich.Spectrum.support,
EnvelopeProperty.Berkovich.Spectrum.preimage_support_zeroLocus,
EnvelopeProperty.Berkovich.Spectrum.continuous_support,
EnvelopeProperty.Berkovich.Spectrum.support_comap}
\leanok
\uses{spectrum:topology, spectrum:pullback, spectrum:quotient-zero-locus,
spectrum:quotient-zero-locus-closed}
For $x\in\mathcal S(A)$ its support is the prime ideal
$\mathfrak p_x=\{a\in A:|a|_x=0\}$. The support map
$\mathcal S(A)\to\Spec A$ is continuous and commutes with restriction of seminorms.
The inverse image of $V(I)$ is the seminorm zero locus of $I$.
\end{definition}
\begin{proof}
The zero set is an ideal by the ultrametric inequality and multiplicativity.
It is proper because $|1|_x=1$; multiplicativity implies that a product has value
zero only if a factor does, proving primality. A support contains $I$ exactly when
every element of $I$ has value zero. This proves the inverse-image formula and
continuity, since such vanishing conditions are closed and all Zariski closed
sets are of the form $V(I)$. Evaluation after restriction is evaluation on the
image of a ring element, proving functoriality of supports.
\end{proof}

\begin{lemma}[Absolute supports are surjective]\label{spectrum:support-surjective}
\lean{EnvelopeProperty.Berkovich.Spectrum.primePoint,
EnvelopeProperty.Berkovich.Spectrum.support_primePoint,
EnvelopeProperty.Berkovich.Spectrum.support_surjective}
\leanok
\uses{spectrum:support, spectrum:quotient-homeomorphism}
Every prime ideal $\mathfrak p$ of $A$ is the support of a real-valued seminorm:
pull back the trivial seminorm of the domain $A/\mathfrak p$.
\end{lemma}
\begin{proof}
The trivial seminorm on a domain assigns zero to zero and one to every other
element. Its pullback to $A$ has value zero precisely on the kernel of
$A\to A/\mathfrak p$, which is $\mathfrak p$. Its support is therefore the
prescribed prime, giving surjectivity of the support map.
\end{proof}

\begin{lemma}[Radicals are detected by absolute seminorms]\label{spectrum:radical-vanishing}
\lean{EnvelopeProperty.Berkovich.Spectrum.zeroLocus_radical,
EnvelopeProperty.Berkovich.Spectrum.zeroLocus_subset_zeroLocus_iff,
EnvelopeProperty.Berkovich.Spectrum.zeroLocus_eq_zeroLocus_iff,
EnvelopeProperty.Berkovich.Spectrum.val_eq_zero_forall_iff_isNilpotent,
EnvelopeProperty.Berkovich.Spectrum.val_eq_of_sub_isNilpotent,
EnvelopeProperty.Berkovich.Spectrum.logAbs_eq_of_sub_isNilpotent,
EnvelopeProperty.Berkovich.Spectrum.zeroLocus_nilradical}
\leanok
\uses{spectrum:support, spectrum:support-surjective}
For ideals $I,J$ of $A$, their seminorm zero loci satisfy
$V(I)^\mathcal S=V(\sqrt I)^\mathcal S$ and
$V(I)^\mathcal S\subseteq V(J)^\mathcal S$ exactly when $J\subseteq\sqrt I$.
They are equal exactly when $\sqrt I=\sqrt J$.
An element has value zero at every point of $\mathcal S(A)$ exactly when it is nilpotent.
Functions differing by a nilpotent have equal seminorms and equal logarithmic evaluations.
The nilradical vanishes at every seminorm.
\end{lemma}
\begin{proof}
Every support is prime, so it contains $I$ exactly when it contains $\sqrt I$.
Because the support map is surjective, inclusion or equality of its inverse images
is equivalent to inclusion or equality of the corresponding Zariski zero loci.
The intersection of the primes containing $I$ is $\sqrt I$, which gives the
inclusion criterion and then the equality criterion. Similarly an element vanishes
at every seminorm precisely when it belongs to every prime, that is, to the
nilradical. Elements of the nilradical are precisely the nilpotents.
If $a-b$ is nilpotent, its seminorm is zero. The ultrametric inequality applied
to $a=b+(a-b)$ and $b=a-(a-b)$ gives $|a|\le|b|$ and $|b|\le|a|$,
respectively; thus their values and logarithms agree.
\end{proof}

\begin{lemma}[Relative supports lie over the prescribed base support]
\label{spectrum:relative-support}
\lean{EnvelopeProperty.Berkovich.RelativeSpectrum.support,
EnvelopeProperty.Berkovich.RelativeSpectrum.continuous_support,
EnvelopeProperty.Berkovich.RelativeSpectrum.support_comap,
EnvelopeProperty.Berkovich.RelativeSpectrum.support_base,
EnvelopeProperty.Berkovich.RelativeSpectrum.supp_comap_algebraMap}
\leanok
\uses{spectrum:relative, spectrum:support}
Let $R$ be a commutative ring, $v$ a real-valued multiplicative ultrametric seminorm
on $R$, and $A$ an $R$-algebra. The support map from $\mathcal S_v(A)$ to $\Spec A$
is continuous and commutes with restriction along $R$-algebra homomorphisms.
For every $x\in\mathcal S_v(A)$,
$\mathfrak p_x\cap R=\mathfrak p_v$, where contraction is taken along $R\to A$.
No surjectivity assertion for this relative support map is made.
\end{lemma}
\begin{proof}
The relative support map is the restriction of the absolute support map to a
subspace, hence is continuous and retains its functoriality. For $r\in R$,
the prescribed restriction gives $|r|_x=|r|_v$. Thus $r$ belongs to the contraction
of $\mathfrak p_x$ exactly when $|r|_v=0$, proving the contraction equality.
\end{proof}
